\begin{frame}[standout]

    Focus sur Neo4j

    \note{
        Neo4j sert d'exemple concret pour montrer comment une base graphe
        répond à des questions de chemin et de voisinage.
    }

\end{frame}

\begin{frame}{Étude de cas}

    https://neo4j.com/customer-stories/

    https://neo4j.com/customer-stories/uber/

    LinkedIn et Neo4j (recommandations de connexions).

    https://neo4j.com/blog/developer/exploring-linkedin-in-neo4j/

    https://www.linkedin.com/posts/rohith-p-561705220_neo4j-graphdatabase-linkedin-activity-7374069650596261889-ClJB

    Problématique.
    Solution NoSQL choisie.
    Architecture mise en place.
    Résultats (performances, coûts).

    \note{
        LinkedIn permet d'illustrer les recommandations de connexions :
        la valeur vient des relations entre personnes, entreprises et intérêts.
        Demander quelles relations et quels parcours seraient nécessaires pour
        proposer une recommandation. Comparer avec le coût de plusieurs jointures
        dans un modèle relationnel.
    }

\end{frame}
